\chapter{The smallest honest Gaussian world}\label{ch:world-schema}
The first active world should be small enough that every represented
quantity has an unambiguous meaning. Small does not mean vague. A two-cell
model with mixed units and hidden replenishment is less interpretable than
a larger model with explicit boundaries.

\section{Physical state and reference are different records}
The stock vector $x$ changes when physical actions or forcing occur. The
reference $x^*$ and scales $\sigma$ are fixed configuration in the first
domain. Both state and reference are nonnegative and share total mass $M$;
every scale is finite and strictly positive. NaN, infinity, a zero scale
or a negative stock is not a small deviation to smooth away.

The potential is
\begin{equation}
 V(x)=\tfrac12\sum_i((x_i-x_i^*)/\sigma_i)^2.
\end{equation}
It is a derived value, not another physical carrier. Storing a cached value
may be convenient, but replay must be able to recompute it from the exact
declared state and configuration. Updating the cache without the state, or
the state without invalidating the cache, risks valuing a nonexistent world.

\section{Accounts belong beside, not inside, the stock vector}
The proposed augmented record includes physical state $x$, persistent
capacity $B$, signed external audit $J$, and random-stream replay coordinates.
They have linked versions but distinct meanings. Adding an EBU balance to a
water total is dimensionally invalid. Treating a physical capacity constraint
as a spendable balance is semantically invalid even if both are numbers.

The old candidate C3 export budget also used a letter resembling $B$.
That does not make it the new persistent account. A safe specification
uses explicit type names and conversions where necessary, not a shared
generic field called budget for every purpose.

\section{Atomic action declarations}
An action records its identity, source, destination, finite quantity, carrier,
write support and relevant constraints. For a lossless transfer,
$\delta_a=q_a(-e_i+e_j)$. A group retains every child and its accepted
increment; it also records the aggregate $\delta_G$.

The source-side quantity convention must remain explicit when importing
historical rate-based examples. A duration belongs in the conversion
$q=\Delta t j$, not in both the conversion and the finite evaluator.
Non-unit-duration fixtures expose errors that unit-duration examples hide.

\section{What is deliberately absent}
The active first domain contains no Allee regeneration, lower/upper
homeostatic band, hidden demand-serving controller or autonomous gradient
flow. Its reference is not an empirically fitted stationary mean by default.
Its scales are not automatically measurement-error bars. It contains one
homogeneous conserved scalar, not an uncalibrated mixture of energy, water,
medicine and welfare.

This is a scope statement, not a claim those phenomena are unimportant.
Adding them later changes state, dynamics or valuation and must receive
its own derivation and registration. An honest small world provides a
controlled question; an underspecified broad world provides ambiguous answers.

\section*{A configuration audit}
Suppose a proposed file gives $x=(18,2)$, reference $(10,10)$ and
scales $(2,2)$. It passes the elementary domain checks. Replace the reference
with $(12,12)$ and closed mass compatibility fails. Replace a scale with
zero and valuation is undefined. Add a regenerative source label without
a regeneration law and the label has supplied no physical mechanism.

Each failure should identify the exact predicate and field. ``Invalid
configuration'' alone is difficult to audit. Silently normalizing the
reference or clipping a stock would make the accepted model differ from
the author's declaration.

\section*{Chapter recap}
Typed state, fixed geometry and complete action declarations make later
checks meaningful. The first world is intentionally limited; its limitations
must remain readable rather than filled with inherited historical defaults.
